\documentclass[11pt]{article}

% Suppress volatile pdfTeX metadata.  The build driver also fixes
% SOURCE_DATE_EPOCH so the companion PDF is reproducible.
\ifdefined\pdfinfoomitdate\pdfinfoomitdate=1\fi
\ifdefined\pdftrailerid\pdftrailerid{}\fi
\ifdefined\pdfsuppressptexinfo\pdfsuppressptexinfo=15\fi

\usepackage[margin=0.92in]{geometry}
\usepackage[numbers,sort&compress]{natbib}
\usepackage{microtype}
\usepackage{xcolor}
\usepackage{hyperref}

\definecolor{linkblue}{HTML}{1F5A7A}
\hypersetup{
  colorlinks=true,
  linkcolor=linkblue,
  citecolor=linkblue,
  urlcolor=linkblue,
  pdftitle={Can local quantum rules reproduce a Dirac particle? A plain-language account},
  pdfauthor={Douglas Clark}
}

\setlength{\parindent}{0pt}
\setlength{\parskip}{0.55em}
\setlength{\emergencystretch}{2em}
\renewcommand{\baselinestretch}{1.04}
\widowpenalty=10000
\clubpenalty=10000
\displaywidowpenalty=10000

% Generated from content-addressed v0.4 ledgers; do not edit by hand.
\newcommand{\ProductOperatorMax}{0.03000726}
\newcommand{\ProductFineOperatorMax}{0.007496690}
\newcommand{\DirectOperatorMax}{0.002759767}
\newcommand{\DirectFineOperatorMax}{0.0006878673}
\newcommand{\DirectIndependentMax}{5.0626\times 10^{-17}}
\newcommand{\ProductDevHash}{\texttt{c9f445358a7ada5d8d20208c485529e6c6da93768f2ee1c9dea76afed2baab7b}}
\newcommand{\ProductConfirmHash}{\texttt{34c16fd959148d99a5e32c4836d69acb33794d548c11c8ac83986990742fac72}}
\newcommand{\DirectDevHash}{\texttt{89d7c65a37bf8d38fea3a87b7f656c02cae57664fcf6b53b725fb8ace7c228c4}}
\newcommand{\DirectConfirmHash}{\texttt{be5a9f85770e1d201abbc269c26da25d821f11f67db364c434473b5c70517e65}}
\newcommand{\ProductPositionTvMax}{0.00194047}
\newcommand{\ProductMomentumTvMax}{2.0061\times 10^{-4}}
\newcommand{\ProductTotalCurrentMax}{0.00844655}
\newcommand{\ProductLocalCurrentMax}{0.00844655}
\newcommand{\ProductInfidelityMax}{4.5070\times 10^{-5}}
\newcommand{\ProductZitterMax}{0.00881467}
\newcommand{\DirectPositionTvMax}{9.8902\times 10^{-5}}
\newcommand{\DirectMomentumTvMax}{1.3971\times 10^{-4}}
\newcommand{\DirectTotalCurrentMax}{3.6737\times 10^{-5}}
\newcommand{\DirectLocalCurrentMax}{1.9639\times 10^{-4}}
\newcommand{\DirectInfidelityMax}{3.1546\times 10^{-7}}
\newcommand{\DirectZitterMax}{2.2840\times 10^{-5}}


\title{\textbf{Can Local Quantum Rules Reproduce a Dirac Particle?}\\[0.35em]
\large A Plain-Language Account of a Prospective, Audited Test}
\author{Douglas Clark}
\date{Lay companion --- July 2026}

\begin{document}

\maketitle

\begin{abstract}
This project asked a deliberately limited question.  If we place quantum
information on a one-dimensional ring and allow it to interact only with
nearby neighbors, can two rulebooks, each fixed before its own counted
outcomes, reproduce the same
carefully chosen part of the behavior of a Dirac particle?  Both registered
rulebooks received positive development and single-use confirmation
classifications, and a separately written high-precision program reconstructed
their recorded computations.  The surviving evidence covers finite time, low
momentum, one particle, declared changes of mass and position-dependent
external potential, and a
fixed collection of accuracy, causality, symmetry, and resource checks.

A later hostile audit materially narrowed the headline.  One empirical monitor
for unwanted zone-edge behavior contained a factor-of-two error in both the
main and replay programs; it selected no modes.  Those empirical values and one
related rough control are withdrawn.  A different, pre-outcome analytic bound
used the correct zone edge and remains far below its frozen threshold.  The
audit also established that the two free one-particle rules are exactly related
by a stroboscopic time-frame change represented by opposite half-cell phases,
after remapping the mass angle.  They are therefore two separately registered specifications, not proof
of two independent architecture classes.  The defensible result is qualified
constructive sufficiency for the intact shared test core.  It is not a
derivation of physics from information, a discovery of the Dirac automaton, or
a model of an interacting quantum field.
\end{abstract}

\section{The question behind the test}

Many physical theories describe change through local rules: what happens here
depends on what is already here and nearby, not on an instantaneous command
from the other side of the universe.  Quantum cellular automata and quantum
walks make that idea unusually explicit.  Space is represented by a grid,
quantum information occupies the grid, and a clock applies a repeating pattern
of small local operations.

This study asks whether such a construction can reproduce a bounded
\emph{operational role} of a Dirac particle.  ``Operational role'' means a
declared bundle of things the construction must do: evolve the right states,
respond correctly to specified interventions, produce the right measured
distributions and currents, respect causal and symmetry conditions, and stay
within fixed error and resource budgets.

The wording matters.  We did not ask whether the construction \emph{is} an
electron, whether the universe is literally a grid, or whether information is
the unique substance of reality.  The target Dirac equation, quantum building
blocks, lattice, clock, constants, input states, measurements, and permitted
interventions were supplied.  What had to be constructed and tested was an
explicit local mechanism under those assumptions.

The target was the one-particle Dirac equation in one space dimension on a
periodic region of length 16, over a time horizon of 2 in the registered units.
Inputs were confined to an 18-dimensional low-momentum sector.  This is a
small, controlled laboratory for a relativistic quantum role, not a claim
about every state, scale, boundary, or duration.

\section{Why matching a curve would not be enough}

Even this reduced Dirac problem combines several demands.  A state has two
coupled directional components.  Its evolution is quantum and reversible.  It
has a maximum causal speed.  Positive- and negative-energy components can
interfere, producing the rapid current oscillation called
\emph{Zitterbewegung}.  Changing the mass or applying a position-dependent
scalar potential should change the evolution in a prescribed way.

A lattice animation can look persuasive while getting some of these
relationships wrong.  We therefore froze a linked collection of tests rather
than judging trajectories by eye.  The primary comparison bounded the largest
possible evolution error across the entire declared input sector, not just a
few attractive example states.  Additional witnesses compared position and
momentum probabilities, total and local current, state fidelity, and the
Zitterbewegung current history.  Other checks covered normalization, causal
support, parity, time reversal, intervention mapping, and counted circuit
resources.

The study also included mass changes and spatial potentials.  These matter
because a construction that works only when left alone has not reproduced the
same operational role as one that responds correctly when the target theory is
intervened on.  The intervention schedules were exact formulas fixed before
their outcomes were opened; they were not fitted site by site after seeing a
result.

There is substantial prior work showing that local quantum walks and cellular
automata can approach Dirac dynamics, including rigorous observational
convergence and the direct one-dimensional Dirac automaton
\citep{Arrighi2014DiracWalk,Bisio2015QuantumFieldQCA}.  Zitterbewegung and
potential scattering in that direct family are also established
\citep{Bisio2013DiracQCADynamics}.  The novelty here is therefore not the
existence or discovery of the automaton.  It is the prospective comparison
protocol: two separately frozen rulebooks, explicit interventions and
resources, single-use confirmation, high-precision reconstruction, and a
claim boundary that records what later audit did and did not preserve.

\section{The two registered rulebooks}

The first rulebook was an \emph{ordered product walk}.  In each small time step
it applied three kinds of local action in a declared order: a phase associated
with the scalar potential, an on-site mixing associated with mass, and
conditional transport that moved the two directional components to neighboring
sites.  Repeating these simple pieces approximates the combined Dirac
evolution.  Because the pieces are applied one after another, their ordering
creates a leading finite-step error.  The registration therefore expected the
certified error to fall by roughly one half whenever the spatial and temporal
resolution doubled.

The second rulebook used the established \emph{direct Dirac quantum cellular
automaton} (QCA).  Its one-particle update was realized locally as an edge operation
followed by an on-site routing swap.  In free evolution its leading global
error was expected to fall by roughly a factor of four under refinement;
splitting in a spatial potential restored a leading factor-of-two pattern.
The direct registration had its own lattice sizes, thresholds, and resource
clock, frozen for that construction rather than copied from the first one.

These descriptions sound architecturally different, and the registered gate
partitions really are different.  But the later algebraic audit found an exact
relationship underneath them.  After translating the mass parameter, the free
one-particle updates are related by a stroboscopic ``time-frame'' change,
represented mathematically by opposite half-cell phases on the two directional
components \citep{AsbothObuse2013TimeFrames}.  Those phases are not themselves
registered integer-cell transport layers.  In this same one-excitation
description, neighboring product gates can also be fused into the direct
parameterized form.  Consequently, the declared
three-layer versus two-layer counts are faithful descriptions of the frozen
protocol clocks, but they are not minimal hardware costs.  The experiment
covers two registered specifications.  It does not establish two algebraically
independent laws or a broad class of inequivalent architectures.

\section{How the test tried to prevent moving the goalposts}

The product study completed first.  The direct study was then registered as a
separate follow-on, so this was not a simultaneous two-arm trial.  Within each
study, however, the scientific rule, target, states, interventions,
measurements, thresholds, numerical precision, and stopping logic were sealed
before the counted outcomes.

Each rulebook first faced one indivisible development batch containing four
scenarios.  Only a committed development pass accepted by the replay could
open a fresh confirmation batch.  Confirmation contained two new intervention
scenarios and could be used only once.  The held-out cases were not intended as
a large statistical sample.  Their purpose was narrower: to check that the
unchanged construction and maps still met the full frozen bundle when given
new exact schedules.

The main dynamics used 256-bit outward-rounded arithmetic.  Rather than
reporting a floating-point estimate and hoping rounding was harmless, the
calculation carried certified enclosures that had to contain the true numerical
answer.  A separately written implementation then rebuilt the candidate and
target evolutions at 110 decimal digits through different numerical routes.
This was an algorithmically separate reconstruction by the same project, not
an external replication.

The audit trail includes failed infrastructure attempts.  One used a
mathematical tool that assumed real rather than complex data.  Another created
part of the calculation at ordinary machine precision before switching to the
required high precision.  A third produced a rigorous error bar that became too
wide under a spatial potential.  Each stopped under its frozen rules before
confirmation.  Those failures show why plausible decimal outputs are weaker
than a successfully reconstructed certified bound.

\section{What the recorded results showed}

All four counted version-0.4 permanent result records (called
``ledgers'')---development and confirmation for each rulebook---recorded
positive classifications.  The product walk covered four
development scenarios and two confirmation scenarios, as did the direct QCA.
The intact shared checks included the uniform evolution bound, observable
witnesses, causal support, symmetry requirements, resource accounting, and
algorithmically separate reconstruction.

For the product walk, the largest certified evolution distances across its
scenarios fell from \ProductOperatorMax{} to about 0.0150 to
\ProductFineOperatorMax{} as resolution doubled.
For the direct rulebook, the largest free or mass-only distances fell from
about 0.000560 to 0.000140 to 0.00003498.  Its potential-split distances fell
from \DirectOperatorMax{} to about 0.00138 to \DirectFineOperatorMax{}.  These
are upper bounds, so smaller is better.  Their refinement patterns matched the separately frozen
expectations.

Those numbers should not be used as a race between the rulebooks.  The two
registrations used different lattice sizes, schedules, gates, and compilation
clocks.  The meaningful result is that each met its own intact, prospectively
justified criteria---excluding the direct empirical zone-edge monitor---and that
the detailed position, momentum, current, fidelity,
Zitterbewegung, causal, and symmetry checks supported the same bounded role
core.

At the finest declared resolutions, the product construction counted 384
serial rounds of local gates for a free run and 448 with an active potential.  The direct
construction counted 320 free layers, 400 for one potential schedule, and 416
for the other.  These totals expose rather than hide the work charged by each
registration.  Because gate fusion can change them, they do not establish a
fundamental efficiency advantage.

\section{The audit that changed the headline}

The direct QCA has a special issue near the far edge of the lattice momentum
zone.  The registered low-momentum inputs were kept far away from that region,
and the protocol included both an analytic leakage bound and an empirical
monitor intended to measure any probability that reached it.

After the result, a hostile-reader audit found a factor-of-two error in the
empirical monitor.  The intended zone-edge condition began at three quarters
of the maximum lattice momentum.  The implemented integer test instead placed
the boundary beyond every momentum that could exist on the lattice.  It
therefore selected an empty set and returned zero regardless of the state.
The main measurement program, the separately written replay, and one extra
stress diagnostic that did not decide pass or fail all copied the same mistaken
condition.

All empirical zone-edge probabilities, their pass flags, and the zero reported
by that stress diagnostic are withdrawn.  They are not repaired in place, and the historical
ledgers are not rewritten.  An audit-only use of the correct condition makes
that rough stress diagnostic nonzero, about $3.49\times10^{-8}$, confirming
that the old zero was vacuous.  This was a post-confirmation, non-gating rough
potential case, not one of the registered analytic leakage schedules described
next.

The analytic leakage calculation was a different route and used the correct
zone-edge boundary.  Its largest pre-outcome bound was about
$1.84\times10^{-164}$ against a frozen threshold of $10^{-24}$.  That analytic
support remains intact.  So the defect does not erase the low-band result, but
it does mean the complete originally registered direct role was not correctly
executed.  The surviving conclusion must carry that qualification.

This episode also exposes an important limit of ``independent'' software.
Different algorithms can still share the same misunderstanding of a written
specification.  Replay was valuable for precision, construction, and numerical
separation, but it did not independently check the written meaning of this
predicate.  Reconstruction, specification review, and hostile audit solve
different problems; a strong workflow needs all three.

\section{What survives, and what does not follow}

The intact evidence supports the following bounded statement:

\begin{quote}
Two separately registered, strictly local quantum update specifications are
constructively sufficient for the intact shared core of a declared
one-dimensional Dirac role over the tested low-momentum, one-particle,
finite-time domain, under the frozen interventions, precision, resources, and
verification rules.  The direct registration's empirical zone-edge monitor is
withdrawn.
\end{quote}

``Constructively sufficient'' means that explicit mechanisms were built and
met the intact stated tests.  ``Qualified'' means that the invalid monitor is excluded
and that the claim stops at the supplied domain.  ``Two specifications'' means
two prospective records with different declared decompositions and numerical
signatures, not two proven architecture classes.

Several stronger conclusions do not follow:

\begin{itemize}
  \item The study does not derive quantum theory, spacetime, the Dirac
  equation, or particles from information alone.
  \item It does not show that the universe is literally a cellular automaton
  or that information is the unique foundation of nature.
  \item It does not discover the direct Dirac automaton, its continuum limit,
  Zitterbewegung, or potential scattering; those are prior results.
  \item It does not establish accuracy for every momentum, duration, mass,
  potential, boundary condition, or resolution.
  \item It does not show that the registered layer and gate counts are minimal
  or predict a hardware advantage.
  \item It does not establish necessity, uniqueness, or architecture
  independence.
\end{itemize}

\section{The spectrum and many-particle boundaries}

The low-momentum restriction is scientifically important.  Near the opposite
edge of the lattice momentum zone, the direct automaton has an exact partner at
high quasienergy, the discrete-time analogue of band energy.  After removing a
global phase, this partner behaves locally
like a Dirac mode with the opposite mass sign.  Recent work distinguishes such
a high-quasienergy \emph{pseudodoubler} from an ordinary extra low-energy
particle \citep{Gupta2026FermionDoubling,Gupta2025DiracVacuum}.  It does not
invalidate the registered low-band comparison, whose valid analytic bound
charges leakage toward that region.  It does rule out saying that the same
construction is Dirac-equivalent across the whole momentum zone or simply
``doubler-free.''

The one-particle boundary is equally important.  The direct local circuit is a
well-defined hard-core qubit automaton---its local subcells do not allow more
than one excitation---but its phase
when two excitations meet is not the phase of the standard many-fermion
extension.  The
product registration did not freeze the corresponding vacuum and
double-occupation phases in a directly comparable way.  Because every primary
test used one excitation, this study establishes no collision, interacting
field, filled-vacuum, particle-creation, or product-versus-direct many-particle
claim.  Such questions require a new, symmetrically specified prospective
experiment; published multiparticle constructions make clear that extra
structure is needed \citep{Farrelly2014CausalFermions,Mlodinow2020QFTQCA}.

\section{Why the result is worth having}

The scientific value is not a surprising new lattice formula.  It is a worked
example of how to ask a construction question without allowing the claim to
expand after the result looks encouraging.

The test made the target behavior executable rather than rhetorical.  It
charged local operations and numerical uncertainty.  It separated development
from single-use confirmation.  It preserved failed verification lineages.  It
reconstructed the calculation through a second algorithm.  Most importantly,
when later analysis found both an exact equivalence and a shared reducer bug,
the project narrowed the claim instead of protecting the original headline.

That combination leaves a result that is smaller but more durable.  The two
rulebooks provide disciplined coverage of two frozen implementations of a
common bounded role.  The exact time-frame relationship explains why that
coverage is not architecture-class breadth.  The valid analytic leakage bound
explains why the low-band result survives.  The withdrawn empirical monitor
shows why cryptographic fingerprints and replay are valuable but not
sufficient.

The appropriate next science is therefore a separate registration: for
example, a genuinely distinct local construction defined under an explicit
equivalence rule, or a symmetric many-particle test with both many-particle rules
fixed in advance.  It should not be added retroactively to make the present
paper look broader.

\section{Conclusion}

Within a finite, low-momentum, one-particle domain, two separately registered
local quantum rulebooks met the intact shared core of a demanding Dirac test.
Their ledgers recorded positive development and single-use confirmation
classifications, and a separately implemented high-precision route reproduced
the committed computations as implemented.  A later audit withdrew an
empirical zone-edge
monitor shared by both programs while preserving a correctly derived analytic
leakage bound.  The same audit showed that the free one-particle rules are
exactly related by a change of time frame, so they cannot support a claim of
independent architecture classes.

The honest conclusion is qualified constructive sufficiency under declared
inputs and boundaries.  That is less dramatic than deriving reality from
information.  It is also more useful: it says exactly what was built, exactly
what was tested, exactly what broke under audit, and exactly what evidence is
left.

\section*{Evidence and reproducibility}

The technical manuscript, qualified result, post-result audit, generated data,
and build scripts are archived together in the project repository.  This
companion contains no figures, and its numerical result macros are generated
from the same content-addressed ledgers used by the technical paper.  The
repository manuscript README records the exact paths and build command.

\bibliographystyle{unsrtnat}
\bibliography{quantum_dirac_paper/bibliography}

\end{document}
